% Manuscript consolidation draft, 2026-10-21.
% Authoritative inputs: the existing architecture, protocol, and experimental
% methods sources in this manuscript branch; core source implementation pinned
% at Git revision 17c7453c50793643b590fb5beac6ad77b9aa53ce.
% Standalone Section 4: do not also include the superseded Sections 5 and 6.
% No new results, references, experiment rows, mathematical labels, or figures.

\section{Materials and Methods}
\label{sec:04}
\label{sec:architecture}

This section presents the research design, implemented cryptographic voting
protocol, and experimental methodology. It first identifies the changes made
to the previously published biometric voting architecture and defines the
notation used to describe the new ballot-processing pathway. It then follows
the sequence of election initialization, candidate-blind biometric
authorization, local encrypted-ballot construction, proof verification,
homomorphic accumulation, and threshold result recovery. The final
subsections document the implemented software, controlled workloads, timing
boundaries, and statistical treatment. Security assumptions and residual
adversarial capabilities have already been established in
Section~\ref{sec:system-model}. The methodology described here concerns the
new cryptographic investigation, not a repeat of the earlier biometric
recognition experiments.

\subsection{Research Design and System Architecture}
\label{subsec:architecture-overview}

The study uses a system-design and controlled experimental evaluation
approach to examine the integration of biometric authorization with
cryptographically validated, encrypted voting. Its starting point is the
previously published \emph{Blockchain\textendash{}integrated IoT biometric
architecture for secure and decentralized electronic voting}
\citep{rifat2026biometric}. That work combined fingerprint and facial
verification, identity-aware application services, and blockchain voting
operations. The present study retains the authentication foundation while
changing how an authorized person prepares, submits, and has a ballot counted.
In particular, candidate selection is excluded from the inspected
biometric-authorization request, ballot admissibility is expressed as a
zero-knowledge relation, and accepted votes contribute encrypted points
rather than plaintext candidate counts during the voting period.

The architecture connects the enrollment and biometric services, credential
and election management services, a voter-facing web client with its local
privacy worker, an Ethereum-compatible voting contract, and the tally
trustees. The biometric service determines which enrolled person is
requesting authorization, while the credential service controls release of
the applicable secret and eligibility witness. Candidate selection,
encryption, and proof generation take place in the client-side privacy
worker. The smart contract verifies the submitted proof and nullifier before
updating the encrypted aggregate. After voting closes, three of five
trustees can contribute partial decryptions, whose consistency is checked
outside the contract before the administrator submits the recovered result.
These are separate verification boundaries rather than a single
identity-to-tally security mechanism.

Two independent commitment structures coexist in the application.
Semaphore-based identity registration and the \texttt{ZKLocalVoting}
contract provide the existing registration and election-anchoring pathway.
The encrypted-voting pathway instead uses a separately generated custom
credential, an election-specific Poseidon commitment tree, and the
\texttt{VoteChainDirect} contract. The validity circuit proves membership
against the custom eligibility root frozen for the relevant election. A
Semaphore group root cannot be substituted for that root. The
\texttt{ZKLocalVoting} contract also exposes a separate candidate-indexed
voting method, which is not part of the encrypted-ballot measurements.
Table~\ref{tab:architecture-prior-comparison} records the architectural
differences. It is a comparison of information flow and contract
responsibilities, not a matched transaction-cost experiment.

\begin{table}[tbp]
\centering
\caption{Voting-path architecture in the previously published system and the present encrypted-ballot extension. The legacy column describes
the earlier published design; the current column describes the implemented
private-ballot pathway.}
\label{tab:architecture-prior-comparison}
\footnotesize
\setlength{\tabcolsep}{4pt}
\renewcommand{\arraystretch}{1.12}
\begin{tabular}{@{}>{\raggedright\arraybackslash}p{0.18\linewidth}>{\raggedright\arraybackslash}p{0.34\linewidth}>{\raggedright\arraybackslash}p{0.40\linewidth}@{}}
\toprule
\textbf{Aspect} & \textbf{Previously published system (2026)} & \textbf{Proposed encrypted-ballot pathway} \\
\midrule
Biometric authorization &
Face/fingerprint verification authorizes voting through identity-aware services. &
Inherited face/fingerprint services authorize release of a protected voting credential; the circuit does not prove biometric matching. \\
\addlinespace[3pt]
Eligibility credential &
Salted pseudonymous voter hash, derived using the registered identity; checked against the contract's voter registry. &
Separate custom Poseidon credential commitment and frozen depth-17 Merkle root; the existing Semaphore registration anchor is not this proof root. \\
\addlinespace[3pt]
Candidate information flow &
The voting application prepares a transaction containing the selected candidate identifier. &
Candidate selection stays in the web client during biometric authorization and enters the local privacy worker for encryption. \\
\addlinespace[3pt]
Ballot submitted on-chain &
Candidate identifier and registered pseudonymous voter hash supplied to \texttt{LocalVoting.castVote}. &
Groth16 proof, election-scoped nullifier, and 16 encrypted candidate-slot ciphertexts supplied to \texttt{VoteChainDirect.castVote}; no plaintext candidate identifier. \\
\addlinespace[3pt]
Ballot acceptance &
Checks registered voter hash, unused vote status, and valid candidate identifier. &
Checks election phase, unused nullifier, canonical ciphertext coordinates, and a 73-public-signal Groth16 ballot-validity statement. \\
\addlinespace[3pt]
Tally during voting &
Plaintext \texttt{voteCount} incremented for the selected candidate. &
BabyJub exponential-ElGamal ciphertexts accumulated for fixed candidate slots; no plaintext per-candidate tally during voting. \\
\addlinespace[3pt]
Result disclosure &
Candidate counts directly readable from the contract; no trustee-decryption step. &
Three-of-five trustee partial decryptions checked off-chain; final totals and transcript hash recorded on-chain after separate state and arithmetic checks. \\
\addlinespace[3pt]
Residual trust &
Identity-aware service can reconstruct the mapping between voter identity and pseudonymous hash. &
Credential custodian can derive an identified voter's nullifier; ballot secrecy further depends on client integrity and trustee-share custody. \\
\bottomrule
\end{tabular}
\end{table}

\subsection{Cryptographic Preliminaries and Notation}
\label{subsec:protocol-notation}
\label{sec:05}
\label{sec:cryptographic-protocol}

The voting relation is represented as a Circom circuit over the BN254
scalar field $\mathbb{F}$. Poseidon is used to compute credential
commitments, membership-tree nodes, election-scoped nullifiers, and
election-context bindings \citep{grassi2021poseidon}. Ballot encryption
uses exponential ElGamal in the configured prime-order BabyJub subgroup
$\mathbb{G}$. Groth16 provides a succinct argument for the combined
eligibility and ballot-construction relation
\citep{groth2016snark}. These primitives have established prior uses.
Their role here is to support one concrete proof statement verified
during encrypted-ballot acceptance, followed by separately checkable
threshold recovery.

Table~\ref{tab:protocol-notation} defines the notation used in the
remaining methodology. The configured circuit has $M=16$ encrypted
candidate slots and a fixed Merkle depth of $L=17$. The active
candidate count $K$ satisfies $1\leq K\leq M$, so changing $K$ does not
change the fixed-width ciphertext representation or the circuit depth.
The Merkle witness supplies one sibling and one direction bit per
tree level. The public input vector contains nine election and nullifier
values followed by four coordinate arrays of length $M$:

\begin{center}
\begin{minipage}{\linewidth}
\centering
\makeatletter\def\@captype{table}\makeatother
\caption{Consolidated notation for the private-voting and tally protocol.}
\label{tab:protocol-notation}
\footnotesize
\begin{tabular}{@{}p{0.29\linewidth}p{0.63\linewidth}@{}}
\toprule
\textbf{Symbol} & \textbf{Meaning} \\
\midrule
$e$, $s$ & Election identifier and private voter credential secret \\
$\mathbb{F}$, $\mathbb{G}$, $q$ & Circuit field, prime-order BabyJub subgroup, and subgroup order \\
$G$, $Y=(Y_x,Y_y)$ & Subgroup generator and election encryption public key \\
$d_{\mathrm{id}}$, $d_{\mathrm{null}}$, $d_{\mathrm{ctx}}$ & Poseidon domain separators for identity, nullifier, and election context \\
$H_P$ & Poseidon hash with the ordered inputs shown in each equation \\
$\rho_e$, $p$ & Frozen election eligibility root and private Merkle authentication path \\
$\mu_e$, $\delta$, $\chi_e$ & Candidate-manifest hash, contract-domain value, and election-context hash \\
$K$, $M$, $L$ & Active candidates, candidate-slot capacity ($16$), and Merkle depth ($17$) \\
$C_s$, $n_e$ & Voter identity commitment and election-scoped nullifier \\
$v_j$, $r_j$ & Private ballot bit and encryption scalar for slot $j$ \\
$C_j=(R_j,S_j)$ & Exponential-ElGamal ciphertext for ballot slot $j$ \\
$\mathbf{w}$, $\mathbf{u}$, $\pi$ & Private circuit witness, public-signal vector, and Groth16 proof \\
$\mathcal{R}_{\mathrm{VC}}$ & Constrained relation verified by the direct voting circuit \\
$B$, $\overline{C}_j$ & Number of accepted ballots and encrypted aggregate for slot $j$ \\
$x_i$, $X_i$ & Trustee $i$'s private share and corresponding public share $x_iG$ \\
$D_{i,j}$, $\gamma_{i,j}$ & Trustee $i$'s partial decryption and its consistency proof for slot $j$ \\
$I$, $\lambda_i$, $t_j$ & Valid trustee set, reconstruction coefficient, and recovered slot total \\
$T_{\mathrm{DKG}}$, $T_{\mathrm{fin}}$ & Distributed-key-generation and finalization transcripts \\
$h_{\mathrm{DKG}}$, $h_{\mathrm{fin}}$ & Canonical SHA-256 hashes of the corresponding transcripts \\
\bottomrule
\end{tabular}
\end{minipage}
\end{center}

\begin{equation}
\label{eq:protocol-public-inputs}
\mathbf{u}=\bigl(
\rho_e,e,\mu_e,K,\delta,Y_x,Y_y,\chi_e,n_e,
\mathbf{R}_x,\mathbf{R}_y,\mathbf{S}_x,\mathbf{S}_y
\bigr).
\end{equation}

The public statement therefore has $9+4M=73$ field signals.
The private witness $\mathbf{w}$ comprises the custom credential secret,
Merkle authentication path, candidate-selection bits, and per-slot
encryption scalars. Neither a biometric measurement nor a civil identifier
is an input to the circuit. The public election-context values are supplied
by the configured contract when it reconstructs the verifier inputs.

\subsection{Voter Registration and Election Initialization}
\label{subsec:architecture-registration}
\label{subsec:protocol-eligibility}

Enrollment takes place through the existing identity-aware registration
service. For an enrolled voter, the application prepares a Semaphore
identity and a separate custom credential with a randomly generated,
nonzero secret $s$. The respective credentials are stored in protected
form alongside the voter record, but support distinct functions.
Semaphore registration does not determine membership in the custom
encrypted-voting relation. For the latter, a Poseidon commitment $C_s$
is calculated with the identity-specific domain separator:

\begin{equation}
\label{eq:protocol-identity}
C_s=H_P(d_{\mathrm{id}},s).
\end{equation}

Before election activation, the election management service gathers
registered custom commitments and constructs a frozen, election-specific
Poseidon Merkle tree of depth $L=17$. A voter proves that the private
credential commitment belongs to this approved set by reconstructing the
root from a private path $p$. At each level, a Boolean direction bit
determines the ordering of the current node and its sibling. The final
root must match the election's configured value $\rho_e$:

\begin{equation}
\label{eq:protocol-merkle}
\operatorname{MerkleRoot}_{L}(C_s,p)=\rho_e.
\end{equation}

This proves membership of a committed credential, not the correctness of
the underlying civil registration or any biometric measurement. The
credential custodian maintains access to the enrolled voting secrets under
the trust assumptions stated in Section~\ref{sec:system-model}.

Election configuration also freezes the ordered candidate manifest and
its hash $\mu_e$. Candidates are ordered by their registered identifiers
and mapped into the first $K$ positions of the 16-slot ballot. The
remaining positions are reserved padding. The administrator additionally
sets the election identifier, context-binding parameters, and collective
encryption public key. Five designated trustees participate in the
distributed-key-generation workflow. Its transcript provides the public
election key and public-share verification material. The election
management service checks that transcript and stores its canonical hash
with the election record.

The election-start procedure deploys or resolves the configured Groth16
verifier and the direct voting contract. It assigns the frozen eligibility
root, election identifier, candidate-manifest hash, candidate count,
contract domain, and encryption key. While the contract remains in its
registration phase, the administrator pins the trustee-transcript hash.
After activation, ballot proofs must satisfy this configured election
statement. Changes to administrative eligibility or candidate ordering
are not silently propagated to an already frozen proof root and manifest.

\subsection{Candidate-Blind Biometric Authorization}
\label{subsec:architecture-biometric-authorization}

The voter-facing application keeps the selected candidate within the
client during biometric authorization. The local privacy worker first
generates an ephemeral public-private key pair to protect credential
delivery. The browser submits the ephemeral public key and the required
authorization data to the face or fingerprint endpoint without transmitting
a plaintext candidate identifier. The inherited biometric service handles
acquisition and recognition through the corresponding interface. Following
successful identification, the credential service resolves the enrolled
record and checks that the election-specific custom voting credential is
available. These checks are identity-aware and remain separate from the
Groth16 voting relation.

The credential service obtains the voter's Merkle witness from the frozen
election snapshot and checks that the reconstructed root matches the
value configured in the direct voting contract. A mismatch prevents
authorization from continuing. It then recovers the stored custom secret
and prepares a protected credential envelope using ephemeral P-256 key
agreement, HKDF-SHA-256, and AES-256-GCM. The response supplies the
encrypted envelope, Merkle witness, election context, and authorization
expiry information. This procedure delivers the credential material
needed for local ballot preparation without passing the intended
candidate through the biometric authorization interface.

The protection is deliberately described as candidate-blind authorization,
not authority-blind anonymity. The biometric and credential services
identify the voter requesting access, and the credential custodian can
recover the stored private secret. Its ability to recompute an
election-scoped nullifier creates the identity-linkage exposure discussed
in Section~\ref{sec:system-model}. Moreover, a valid on-chain ballot
proof demonstrates the prescribed cryptographic relation. It does not
certify that a fresh biometric match was performed at the instant of
transaction submission.

\subsection{Ballot Encoding, Encryption and Nullifier Construction}
\label{subsec:architecture-ballot-submission}
\label{subsec:protocol-encryption}

After authorization, the privacy worker decrypts the temporary credential
envelope using its ephemeral private key. It reconstructs the corresponding
custom commitment and checks the supplied membership path against the
election's eligibility root. The worker also compares the election
identifier and ordered candidate manifest against the public configuration.
The election-context hash $\chi_e$ binds the election identifier, manifest
hash, active candidate count, contract-domain value, and encryption
public-key coordinates:

\begin{equation}
\label{eq:protocol-context}
\chi_e=H_P\bigl(
 d_{\mathrm{ctx}},e,\mu_e,K,\delta,Y_x,Y_y
\bigr).
\end{equation}

The circuit recomputes this value from the public inputs, while the client
checks the candidate ordering against the configured manifest. The circuit
does not reconstruct the entire candidate list from individual candidate
records. These checks establish consistency with the published election
configuration and not the legitimacy of its administrative selection.

The selected candidate is encoded as a private one-hot vector
$\mathbf{v}$ of length $M=16$. Each $v_j$ is Boolean, exactly one
entry is selected, and all slots outside the active candidate range are
constrained to zero:

\begin{equation}
\label{eq:protocol-onehot}
\begin{aligned}
 v_j&\in\{0,1\} && (0\leq j<M),\\
 \sum_{j=0}^{M-1}v_j&=1,\\
 v_j&=0 && (K\leq j<M).
\end{aligned}
\end{equation}

The fixed representation applies to every evaluated active candidate
count. A proof for this relation excludes empty, multiple-choice, and
out-of-manifest ballot vectors under the implemented single-choice
election model. Such constrained encrypted-ballot validity has prior
cryptographic precedent \citep{huber2025ballotvalidity}; the present
method incorporates it alongside custom membership and election
binding.

For each slot $j$, the worker samples a private scalar $r_j$ with
$1\leq r_j<q$ and encrypts the bit $v_j$ under the collective
BabyJub public key $Y$ using exponential ElGamal. With generator $G$,
the ciphertext $C_j$ consists of two curve points:

\begin{equation}
\label{eq:protocol-elgamal}
\begin{aligned}
 R_j &= r_jG,\\
 S_j &= r_jY+v_jG,\\
 C_j &= (R_j,S_j),\qquad 0\leq j<M.
\end{aligned}
\end{equation}

All 16 slots are encrypted, including inactive padding slots. The
voter's chosen candidate and encryption scalars remain in the proof
witness. The circuit checks the encryption-scalar range, supported
curve and subgroup conditions, and ciphertext equations, but does not
establish the quality of the client's random-number generator. The direct
circuit wrapper equates its calculated ciphertexts to the 64 public
coordinate signals. Changing a submitted ciphertext without changing
the corresponding valid witness would therefore invalidate the proof
under the soundness assumption.

For duplicate-ballot prevention, the worker computes the
election-scoped nullifier defined in
Eq.~\eqref{eq:system-nullifier}. The same custom secret and election
identifier yield the same nullifier irrespective of the transaction
signer. The election-specific derivation permits reuse of an eligible
credential in a different configured election without simply repeating
the same public nullifier. It does not prevent one person from voting
through multiple independently issued credentials when enrollment
controls fail.

\subsection{Zero-Knowledge Voting Relation and Proof Generation}
\label{subsec:protocol-verification}

The configured Groth16 relation combines custom commitment membership,
nullifier consistency, election-context binding, a valid one-hot selection,
and correctness of every encrypted slot. Using the public signal vector
$\mathbf{u}$ and private witness $\mathbf{w}$, the relation is expressed as

\begin{equation}
\label{eq:protocol-relation}
\mathcal{R}_{\mathrm{VC}}(\mathbf{u},\mathbf{w})=1
\end{equation}

when the constraints in
Eqs.~\eqref{eq:protocol-identity}--\eqref{eq:protocol-elgamal}
hold for the public values of
Eq.~\eqref{eq:protocol-public-inputs}. The Merkle siblings,
path directions, credential secret, selected candidate, and encryption
randomness remain private. The public proof inputs contain the
election-scoped nullifier, frozen election configuration, and ciphertext
coordinates. The statement does not disclose which active candidate
corresponds to the selected one-hot position.

The client uses the configured circuit and proving key to construct
the Groth16 witness and proof. Before sending a transaction, the worker
compares the resulting 73 public signals against the expected election
configuration, nullifier, and ciphertext array, then performs local
proof verification. A mismatch prevents submission. These local checks
are useful defenses against inconsistent client-side artifacts, but do not
replace contract verification. Proof correctness remains conditional on
the deployed circuit, the associated proving and verification artifacts,
and the cryptographic assumptions of Groth16. The research setup is
development-stage and has not undergone an independently audited
production ceremony.

\subsection{On-Chain Verification and Encrypted Aggregation}
\label{subsec:protocol-onchain-acceptance}

The client transmits the verified proof, election-scoped nullifier, and
64 ciphertext coordinates through the wallet transaction interface to
\texttt{VoteChainDirect.castVote}. No plaintext candidate identifier is
submitted in this call. The contract requires an active election,
checks that the nullifier has not been used, and validates the canonical
field representation of the ciphertext coordinates. It reconstructs
the complete 73-field public statement using the frozen election
configuration and supplied ciphertexts, then invokes the Groth16
verifier. Only a transaction satisfying these checks can mark the
nullifier as used and modify the accepted encrypted tally.

The contract maintains one encrypted aggregate for each of the 16
candidate slots. Every aggregate begins at the BabyJub identity
points, and an accepted ballot is added componentwise to the current
ciphertext state. For $B$ accepted ballots, the encrypted sum in slot
$j$ is

\begin{equation}
\label{eq:protocol-aggregation}
\overline{C}_j=(\overline{R}_j,\overline{S}_j)
 =\sum_{b=1}^{B}C_{b,j}.
\end{equation}

The aggregate represents an encryption of the accumulated candidate
count $t_j$ without making that count plaintext while voting is active.
The contract increments the accepted-ballot count and emits an
acceptance event recording the public nullifier, a ciphertext hash, and
updated count. Public observers can inspect these records and submitted
ciphertexts, but the contract does not expose the voter's plaintext
choice as part of the accepted ballot. Confidentiality additionally
depends on client integrity, encryption security, and threshold-share
custody, as established in Section~\ref{sec:system-model}.

\subsection{Threshold Decryption and Auditable Finalization}
\label{subsec:architecture-finalization}
\label{subsec:protocol-threshold}
\label{subsec:protocol-finalization}

Multi-authority encrypted election tallying has established cryptographic precedent
\citep{cramer1997multiauthority}. The implementation here uses a
three-of-five threshold, with the specific verification checks stated below.

After voting ends, the administrator advances the contract from
\texttt{Active} to \texttt{Ended}. The finalization service retrieves
the accepted-ballot count and 16 encrypted aggregate ciphertexts from
the contract. The implemented recovery procedure uses five trustees
with a reconstruction threshold of three. The previously checked
distributed-key-generation transcript specifies the collective public
key and trustee public shares. Trustee $i$ holds a secret share $x_i$
with corresponding public point $X_i=x_iG$. For candidate slot $j$,
its partial decryption of the aggregate first point is

\begin{equation}
\label{eq:protocol-partial}
D_{i,j}=x_i\overline{R}_j.
\end{equation}

Each partial-decryption record includes a non-interactive
discrete-logarithm equality proof $\gamma_{i,j}$. The verifier must
establish that the same secret share relates the generator $G$ to the
public share $X_i$ and the aggregate point $\overline{R}_j$ to the
submitted partial decryption. The trustee constructs commitments
$A_{1,i,j}=k_{i,j}G$ and
$A_{2,i,j}=k_{i,j}\overline{R}_j$, derives a domain-separated Poseidon
challenge $c_{i,j}$ from the election, trustee, slot, and curve-point
context, and returns
$z_{i,j}=k_{i,j}+c_{i,j}x_i\pmod q$.
The external transcript verifier recomputes the challenge and checks

\begin{equation}
\label{eq:protocol-dleq}
\begin{aligned}
 z_{i,j}G&=A_{1,i,j}+c_{i,j}X_i,\\
 z_{i,j}\overline{R}_j&=A_{2,i,j}+c_{i,j}D_{i,j}.
\end{aligned}
\end{equation}

These equations validate consistency with the public trustee share
and specified encrypted aggregate under the proof assumptions. They
are not executed by the Solidity finalization function.

For at least three distinct trustees with valid evidence, the verifier
combines partial decryptions with Lagrange interpolation coefficients
evaluated at zero:

\begin{equation}
\label{eq:protocol-lagrange}
\begin{aligned}
 \lambda_i&=\prod_{\substack{k\in I\\k\neq i}}
 \frac{-k}{i-k}\pmod q,\\
 D_j&=\sum_{i\in I}\lambda_iD_{i,j}.
\end{aligned}
\end{equation}

It then removes the reconstructed masking term from the aggregate
second point, obtaining an encoded candidate count bounded by the
number of accepted ballots:

\begin{equation}
\label{eq:protocol-tally}
 t_jG=\overline{S}_j-D_j,\qquad 0\leq t_j\leq B.
\end{equation}

The bounded discrete logarithm is recovered by the implemented
baby-step/giant-step procedure. A transcript is rejected if a slot
lacks three valid partial contributions, a padded slot decrypts to a
nonzero count, or active-slot totals do not sum to the accepted-ballot
count. The ability of a threshold coalition to recover aggregate totals
also implies that colluding trustees with sufficient usable shares may
decrypt individually published ballot ciphertexts. The prototype's
locally managed trustees do not establish independently governed
organizational share custody.

The evidence used for independent checking includes the
distributed-key-generation record, the encrypted aggregates,
partial-decryption records and their proofs, and the reconstructed
totals. The application checks the election identifiers, frozen
configuration, ballot count, and supplied aggregates against the
contract's public state. The canonicalized DKG and finalization
transcripts are committed using SHA-256:

\begin{equation}
\label{eq:protocol-transcript-hash}
\begin{aligned}
 h_{\mathrm{DKG}}&=\mathsf{SHA256}
 (\mathsf{Canon}(T_{\mathrm{DKG}})),\\
 h_{\mathrm{fin}}&=\mathsf{SHA256}
 (\mathsf{Canon}(T_{\mathrm{fin}})).
\end{aligned}
\end{equation}

The DKG transcript hash is pinned during election configuration.
Before finalization, the application validates the supplied transcript
against that commitment, verifies trustee contributions, and checks the
aggregate against the ledger. Once the external checks succeed, the
administrator submits the recovered totals and the finalization
transcript hash through \texttt{finalizeTally}. The contract requires the
\texttt{Ended} phase and a previously pinned DKG commitment. It also
requires a nonzero finalization hash, zero totals in padded slots,
active totals no greater than $B$, and a sum of active totals equal to
$B$. If these conditions hold, it stores the plaintext totals and hash
and advances to \texttt{Finalized}.

The on-chain transition does not verify the DKG record, DLEQ proofs,
or derivation of the submitted plaintext counts from the encrypted
aggregate. A committed transcript permits mismatch detection only when
the full evidence can be obtained and independently checked. Therefore,
successful finalization and a matching hash must not be presented as
on-chain proof of correct threshold decryption. On-chain Groth16
verification at ballot acceptance and external verification at result
disclosure provide different, complementary checks.

\subsection{Prototype Implementation}
\label{sec:06}
\label{sec:implementation-experiments}
\label{subsec:implementation-profile}

The prototype combines the inherited biometric services with a
React and TypeScript voter-facing web client and Node.js election
and credential services. The browser privacy worker performs
credential-envelope recovery, election-context checks, ballot
encryption, witness generation, Groth16 proving, and local proof
verification. The private voting relation is compiled from Circom 2
using the fixed 16-slot, depth-17 configuration. The proof workflow
uses \texttt{snarkjs}, while a generated Solidity verifier validates
the 73-signal statement within the
\texttt{VoteChainDirect} contract. The contract enforces election
state and nullifier conditions and maintains the encrypted
BabyJub tally.

The separate threshold implementation prepares and validates the DKG
record, trustee partial decryptions, consistency proofs, aggregate
recovery, and finalization transcript. Its cryptographic checks occur
outside Solidity. The final contract transaction records the result
values and transcript commitment only after its own state and arithmetic
conditions hold. These components form the evaluated encrypted-voting
pathway. The predecessor's fingerprint and facial recognition
measurements were not repeated as part of the four quantitative
campaigns described below.

The cryptographic proving parameters were generated using a
development-stage setup. The experimental transactions were executed
on local Anvil rather than a public blockchain or a permissioned
multi-validator deployment. The reported observations consequently
characterize this implemented research configuration, not production
election throughput, network-finality behavior, or secure distributed
operation of independent trustee institutions.

\subsection{Experimental Procedure and Measurement Scope}
\label{subsec:experimental-environment}
\label{subsec:experimental-campaigns}

The controlled experiments used a Windows-based AMD64 computer
with 16 logical processors and approximately 10.6~GB of available
physical memory. The recorded software environment included
Python~3.10.11, Node.js, Circom, \texttt{snarkjs}, Foundry, and
Anvil. The transaction campaigns used a local Anvil network and
recorded gas consumed by accepted encrypted-ballot calls.
Submission-to-receipt latency was measured from transaction
submission until receipt retrieval. This interval excludes
biometric acquisition, identity authorization, protected credential
delivery, and local ballot proof preparation. Accordingly, it
must not be interpreted as complete voter-session time or as a
public-network confirmation measurement.

Cryptographic measurements were collected independently from
blockchain receipts. The eligible-set campaign timed identity-leaf
preparation, construction of a Merkle tree and target path, witness
generation, Groth16 proof generation, and local proof verification.
The four finalized campaigns were accepted-ballot transactions,
fresh-chain first-ballot execution, active candidate-count
sensitivity, and eligible-set sensitivity. Each addresses a
different computational condition. Table~\ref{tab:implementation-campaigns}
records the measured sample sizes and configuration ranges.

\begin{table}[t]
\centering
\caption{Experimental configurations and measured sample sizes.}
\label{tab:implementation-campaigns}
\small
\begin{tabular}{@{}p{0.23\linewidth}p{0.44\linewidth}p{0.25\linewidth}@{}}
\toprule
\textbf{Campaign} &
\textbf{Experimental configuration} &
\textbf{Measured observations} \\
\midrule
Accepted ballots
& 25 independent elections, with five elections at each
size of 25, 50, 75, 100, and 150 accepted ballots
& 2,000 transactions \\
\addlinespace
Fresh-chain first ballot
& First accepted private ballot following a fresh
local-chain setup
& 20 trials \\
\addlinespace
Candidate count
& Active candidate counts of 2, 4, 8, and 16
with a fixed 16-slot representation
& 20 trials per count \\
\addlinespace
Eligible set
& Eligible populations of 10, 100, 1,000, and
10,000 with Merkle depth fixed at 17
& 20 trials per population \\
\bottomrule
\end{tabular}
\end{table}

The accepted-ballot campaign measured 2,000 successful
transactions across 25 independent elections. Five elections
were run at each target size of 25, 50, 75, 100, and 150
accepted ballots, producing the recorded total of 2,000
observations. An additional 400 warm-up transactions were
excluded from the final statistics. Each accepted transaction
contributed a gas observation and a local submission-to-receipt
interval. Because successive ballots share the evolving state
of an election, this campaign includes both initial and later
accepted ballots rather than characterizing only an empty tally.

The fresh-chain campaign independently isolated the first
accepted encrypted ballot after fresh local-chain initialization.
It recorded 20 measured trials after five excluded warm-ups.
Contract preparation and checks performed after receipt retrieval
lay outside the timed interval. These observations were
analyzed separately from the 2,000-ballot campaign, since the
state and execution conditions define a different measurement
population.

The candidate-count campaign used 2, 4, 8, and 16 active
candidates, with 20 measured transactions per configuration
after five excluded warm-ups. Candidate positions were balanced
within the recorded trial sequence. All trials used the same
16-slot ciphertext representation and compiled proof relation.
Thus, the campaign measures changes in active candidate
configuration within a fixed-width circuit, not a family of
circuits with increasing capacity.

The eligible-set campaign used synthetic populations of
10, 100, 1,000, and 10,000 registered custom commitments.
Each configuration contributed 20 measured runs following
five warm-ups, with execution order varied across sizes.
The selected voter occupied the first indexed position in
the respective fixture. Every run used Merkle depth 17 and
16 encrypted candidate slots. The nominal depth-17 capacity
is 131,072 leaves, but the campaign neither populated that
full capacity nor varied circuit depth. No blockchain
transactions were submitted during this campaign. Its
verification observations represent local proof checking,
not the Solidity verification gas or execution latency.

\subsection{Evaluation Metrics and Statistical Analysis}
\label{subsec:experimental-analysis}

The transaction campaigns report gas consumption and local
submission-to-receipt latency, while the eligible-set campaign
reports individual preparation and cryptographic processing
intervals. The four measurement campaigns retain their own
observation boundaries and are not combined into a fabricated
end-to-end voting time. For every finalized measurement group,
descriptive summaries include sample size, mean, median,
minimum, maximum, sample standard deviation, and
95th percentile. The sample standard deviation is calculated
using the $n-1$ denominator, and the 95th percentile follows
the nearest-rank method. All recorded warm-ups, pilot runs,
and aborted campaigns were excluded as specified by the
experimental procedure. No measured outliers were removed.

The candidate-count experiment is assessed descriptively,
without inferential regression across its four fixed-width
configurations. In the eligible-set experiment, the population
may affect leaf materialization and tree construction, but
the membership circuit continues to constrain a path of
constant depth. The recorded processing stages are therefore
reported individually, with no unmeasured inference about
larger circuits, network-wide throughput, or complete
biometric-to-ballot execution. Proof-generation intervals
reflect the measured software invocation and runtime
boundary rather than isolated primitive timings.

Functional and integration checks supplement the positive
performance campaigns. They cover witness and public-signal
consistency, local Groth16 verification, selected malformed-ballot
and reused-nullifier rejection paths, and threshold-transcript
verification conditions. Such checks do not amount to a
comprehensive adversarial security audit. In particular, the
four quantitative campaigns do not produce a repeated timing
estimate for the complete threshold-finalization workflow.
That stage is therefore described by its implementation and
functional verification evidence, without assigning it an
unsupported end-to-end latency. The corresponding measured
results are presented in the following section.
